\documentclass{article}

\usepackage{fullpage}
\usepackage{color}
\usepackage{amsmath}
\usepackage{url}
\usepackage{verbatim}
\usepackage{graphicx}
\usepackage{parskip}
\usepackage{amssymb}
\usepackage{nicefrac}
\usepackage{listings} % For displaying code
\usepackage{algorithm2e} % pseudo-code
\usepackage[usenames,dvipsnames]{xcolor}

%%
%% Julia definition (c) 2014 Jubobs
%%
\lstdefinelanguage{Julia}%
  {morekeywords={abstract,break,case,catch,const,continue,do,else,elseif,%
      end,export,false,for,function,immutable,import,importall,if,in,%
      macro,module,otherwise,quote,return,switch,true,try,type,typealias,%
      using,while},%
   sensitive=true,%
   alsoother={\$},%
   morecomment=[l]\#,%
   morecomment=[n]{\#=}{=\#},%
   morestring=[s]{"}{"},%
   morestring=[m]{'}{'},%
}[keywords,comments,strings]%

\lstset{%
    language         = Julia,
    basicstyle       = \ttfamily,
    keywordstyle     = \bfseries\color{blue},
    stringstyle      = \color{magenta},
    commentstyle     = \color{ForestGreen},
    showstringspaces = false,
}

% Answers
\def\ans#1{\par\gre{Answer: #1}}
%\def\ans#1{} % Comment this line to produce document with answers

% Colors
\definecolor{blu}{rgb}{0,0,1}
\def\blu#1{{\color{blu}#1}}
\definecolor{gre}{rgb}{0,.5,0}
\def\gre#1{{\color{gre}#1}}
\definecolor{red}{rgb}{1,0,0}
\def\red#1{{\color{red}#1}}
\def\norm#1{\|#1\|}

% Math
\def\R{\mathbb{R}}
\def\argmax{\mathop{\rm arg\,max}}
\def\argmin{\mathop{\rm arg\,min}}
\newcommand{\mat}[1]{\begin{bmatrix}#1\end{bmatrix}}
\newcommand{\alignStar}[1]{\begin{align*}#1\end{align*}}
\def\half{\frac 1 2}

% LaTeX
\newcommand{\fig}[2]{\includegraphics[width=#1\textwidth]{a1/#2}}
\newcommand{\centerfig}[2]{\begin{center}\includegraphics[width=#1\textwidth]{a1/#2}\end{center}}
\def\items#1{\begin{itemize}#1\end{itemize}}
\def\enum#1{\begin{enumerate}#1\end{enumerate}}


\begin{document}

\title{CPSC 440/540 Machine Learning (January-April, 2022)\\Assignment 1 (due Friday January 21st at midnight)}
\author{}
\date{}
\maketitle
\vspace{-4em}

\textbf{IMPORTANT!!!!! Please carefully read the submission instructions that will be posted on Piazza. We will deduct 50\% on assignments that do not follow the instructions}.

Most of the questions below are related to topics covered in CPSC 340, or other courses listed on the prerequisite form. There are several ``notes'' available on the webpage which can help with some relevant background.

If you find this assignment to be difficult overall, that is an early warning sign that you may not be prepared to take CPSC 440
at this time. Future assignments may be longer and more difficult than this one.

We use \blu{blue} to highlight the deliverables that you must answer/do/submit with the assignment.

\section*{Basic Information}


\blu{\enum{
\item Name: Wira Azmoon Ahmad
\item Student ID: 55246318
}}

\pagebreak


\section{Very-Short Answer Questions}



Give a short and concise 1-2 sentence answer to the below questions. All acronyms, symbols, and methods are as used in CPSC 340.

\blu{
\enum{
\item Suppose that a famous machine learing personality is advertising their  ``extremely-deep convolutional fuzzy-genetic Hilbert-long-short adversarial-recurrent neural network'' classifier, which has 500 hyper-parameters. This person claims that if you take 10 different famous datasets, and tune the 500 hyper-parameters based on each dataset's validation set, that you can beat the current best-known validation set error on all 10 datasets. Explain whether or not this amazing claim is likely to be meaningful.
\ans{
It is not meaningful, breaking the golden rule of machine learning, and unlikely to generalise to an actual test set.
}
\item Consider fitting a neural network with one hidden layer. Would we call this a parameteric or a non-parametric model if the hidden layer had $n$ units (where $n$ is the number of training examples)? Would it be parametric or non-parametric if it had $d^2$ units (where $d$ is the number of input features).
\ans{
With $n$ units, it is non-parametric.
With $d^2$ units, it is parametric.
}
\item  Suppose we want to solve the binary classification problem of determining whether 26 by 26 images were images of digits or images of letters. What is an easy way to make our classifier invariant to small translations of the training images?
\ans{
Max pooling can be used, so slight translations would not affect most of the pooled values before classifying.
}
\item Ensemble methods are models that combine the predictions of multiple individual models. Give a reason why an ensemble method could do better than the best individual model in the ensemble.
\ans{
The probability that the ensemble method is correct can be greater than that of any individual model.
}
\item A common method for outlier detection is to collect a large number of outliers, and build a binary classifier that distinguishes these outliers from standard ``inlier'' datapoints. What is an advantage and a drawback of this approach to outlier detection, compared to unsupervised outlier detection methods?
\ans{
An advantage is that complicated outlier patterns can be found,
but the drawback is new types of outliers may not be detected.
}
\pagebreak
\item Why do we typically add a column of $1$ values to $X$ when we do linear regression?
\ans{
This corresponds to the bias variable for each example.
}
\item If a function is convex, what does that say about stationary points of the function? Does convexity imply that a stationary points exists?
\ans{
Convexity does not imply whether or not stationary points exist.
$x^2$ and $|x|$ are convex but the former has a stationary point
while the latter does not.
}
\item What is the cost in $O()$ notation for fitting a linear regression model with $n$ examples and 1 feature using least squares? Explain under what conditions it would make sense to use gradient descent instead of the normal equations for fitting a linear regression model with 1 feature.
\ans{
$O(n)$. With 1 feature, there is no advantage to using gradient descent.
}
\item Suppose you are hired by company to look at a dataset, and they want to ``find which features are relevant'' for their prediction task. What are 2 warnings you would give them regarding the features found by feature selection methods?
\ans{
Some features could have interactions, and are not independent. Irrelevant features could also help by chance.
}
\pagebreak
\item If we fit a linear regression model with L1-regularization of the parameters, what is the effect of the regularization parameter $\lambda$ on the sparsity level of the solution? What is the effect of $\lambda$ on the two parts of the fundamental trade-off (the training error and the amount of overfitting)?
\ans{
An increase in $\lambda$ would correspond to a higher sparsity level of the solution.
The training error would also increase, while the amount of overfitting would decrease.
}
\item Suppose we fit a one-feature linear regression model with a polynomial basis, and this leads to non-zero regression weights. What would the prediction of this model be as the feature value $x^i$ goes to $\infty$? How would this change if you used Guassian RBFs as features?
\ans{
Assuming positive weights, the prediction would tend to $\infty$.
With Gaussian RBFs, this would instead be 0.
}
\item When searching for a good $w$ for a linear classifier, why do we use the logistic loss instead of just minimizing the number of classification errors? 
\ans{
Since the output $y_i$ is numerical, it does not make sense to count the discrete classification errors.
}
\item With stochastic gradient descent, the loss might go up or down each time the parameters are updated. However, we don't actually know which of these cases occurred. Explain why it doesn't make sense to check whether the loss went up/down after each update.
\ans{
Only the expected loss over all examples is expected to go down. It's not meaningful to check after each single update.
}
\pagebreak
\item Suppose we are given two vectors of integers, x=$(x_1,x_2,\dots,x_n)$ and y=$(y_1,y_2,\dots,y_n)$, and we want to find the longest sequence of numbers that appears consecutively in both $x$ and $y$. Describe an $O(n^2)$ algorithm to do this using dynamic programming (you can assume the numbers have a fixed-sized representation).
\ans{
Keep track of an $n\times n$ table $A$, where element $A[i,j]$ is the length of the longest consecutive sequence for $X[1..i],Y[1..j]$. Since for $i,j>1, x_i==x_j, L[i,j]=L[i-1,j-1]+1$, the table can be filled up with DP in $O(n^2)$ time, and the longest sequence can be determined by $max(L)$ and the corresponding indices. Referenced Wikipedia: \url{https://en.wikipedia.org/wiki/Longest_common_substring_problem} .
}
\item Given a vector of $n$ positive integers $(x_1,x_2,\dots,x_n)$, give an $O(n)$ time algorithm for computing, for each position $i$, the sum of all numbers except $x_i$. Next, describe how you could solve this problem in $O(n)$ even if you were not allowed to use subtraction/negation (hint: you can start at the beginning or the end of the list).
\ans{
Store the total sum. For each number, subtract the number itself from the sum.
Without subtraction, use 2 passes. Going through the vector in the forward pass, keep at each index the sum of all smaller indices. Then in the backward pass, add the sum of larger indices. Referenced StackOverflow: \url{https://stackoverflow.com/a/55995896}
}
\item{Consider using a fully-connected neural network for 3-class classification on a problem with $d=10$. If the network has one hidden layer of size $k=100$, how many parameters (including biases), does the network have?}
\ans{
\[(10+1)\times 100 + (100+1)\times 3 = 1403\]
parameters
}
}
}
\pagebreak

\section{Coding Questions}

The questions below use code contained in \emph{a1.zip}, which you can download from the course webpage. The scripts mentioned in the question can be run in Julia's REPL in the directory containing the extracted files. If you have not previously used Julia, there is a list of useful Julia commands (and syntax) among the list of notes on the course webpage.

Several scripts use packages. For example, the \emph{JLD} package is used to load  data and the \emph{PyPlot} package to make  plots. If you have not previously used these packages, they can be installed using:
\begin{verbatim}
using Pkg
Pkg.add("JLD")
Pkg.add("PyPlot")
\end{verbatim}


 
\subsection{K-Means Clustering}

If you run the function \emph{example\_Kmeans}, it will load a dataset with two features and a very obvious clustering structure. It will then apply the $k$-means algorithm with a random initialization. The result of applying the algorithm will thus depend on the randomization, but a typical run might look like this:\\
%\centerfig{.5}{kmeans.png}
(Note that the colours are arbitrary due to the label switching problem.)
But the ``correct'' clustering (that was used to make the data) is something more like this:\\
%\centerfig{.5}{kmeans2.png}

If you run the demo several times, it will find different clusterings. To select among clusterings for a \emph{fixed} value of $k$, one strategy is to minimize the sum of squared distances between examples $x^i$ and their means $w_{y^i}$,
\[
f(w_1,w_2,\dots,w_k,y^1,y^2,\dots,y^n) = \sum_{i=1}^n \norm{x^i - w_{y^i}}_2^2 = \sum_{i=1}^n \sum_{j=1}^d (x_j^i -  w_{y^ij})^2.
\]
 where $y^i$ is the index of the closest mean to $x^i$. This is a natural criterion because the steps of k-means alternately optimize this objective function in terms of the $w_c$ and the $y^i$ values.
 
 \blu{\enum{
 \item Write a new function called \emph{kMeansError} that takes in a dataset $X$, a set of cluster assignments $y$, and a set of cluster means $W$, and computes this objective function. Hand in your code.
 \ans{
 Code in \texttt{kMeansError.jl}.
 \lstinputlisting{a1/kMeansError.jl}
 }
 \item Instead of printing the number of labels that change on each iteration, what trend do you observe if you print the value of \emph{kMeansError} after each iteration of the k-means algorithm?
 \ans{
 The error decreases at each iteration.
 }
 \item Using the \emph{clustering2Dplot} file, output the clustering obtained by running k-means 50 times (with $k=4$) and taking the one with the lowest error. Note that the k-means training function will run much faster if you set \texttt{doPlot = false} or just remove this argument.
 \centerfig{.5}{kmeans50.png}
 \item Explain why the \emph{kMeansError} function should not be used to choose $k$.
 \ans{
 The degenerate case exists where $k=n$ clusters centered at each training example would give 0 error.
 }
 \item Explain why even evaluating the \emph{kMeansError} function on test data still wouldn't be a suitable approach to choosing $k$.
 \ans{
 The model could end up overfitting the test data.
 }
 \pagebreak
 \item The data in \emph{clusterData2.jld} is the exact same as \emph{clusterData.jld}, except it has 4 outliers that are far away from the data. Using the \emph{clustering2Dplot} function, output the clustering obtained by running k-means 50 times (with $k=4$) on \emph{clusterData2.jld} and taking the one with the lowest error. Are you satisfied with the result?
 \ans{
 \centerfig{.5}{kmeans250.png}
 With outliers, the result is not satisfactory especially for the top two clusters.
 }
 \item Implement the $k$-\emph{medians} algorithm, which assigns examples to the nearest $w_c$ in the L1-norm and then updates the $w_c$ by setting them to the ``median" of the points assigned to the cluster (we define the $d$-dimensional median as the concatenation of the medians along each dimension). For this algorithm it makes sense to use the L1-norm version of the error (where $y_i$ now represents the closest median in the L1-norm),
\[
f(w_1,w_2,\dots,w_k,y_1,y_2,\dots,y_n) = \sum_{i=1}^n \norm{x_i - w_{y_i}}_1 = \sum_{i=1}^n \sum_{j=1}^d |x_{ij} - w_{y_ij}|,
\] 
 Hand in your code and plot obtained by taking the clustering with the lowest L1-norm after using  50 random initializations for $k = 4$.
 \ans{
 Code in \texttt{kMedians.jl}.
 \lstinputlisting{a1/kMedians.jl}
 \centerfig{.5}{kmeansl1.png}
 }
}
 }
\pagebreak

\subsection{Regularization and Hyper-Parameter Tuning}

If you run the script \emph{example\_nonLinear} (from Julia's REPL), it will:
\enum{
\item Load a one-dimensional regression dataset.
\item Fit a least-squares linear regression model.
\item Draw a figure showing the training/testing data and what the model looks like.
}
Unfortunately, this is not a great model of the data, and the figure shows that a linear model with a y-intercept of 0 is probably not suitable.
\enum{
\item Update the \emph{leastSquares} function so that it includes a \emph{y-intercept} variable $w_0$, and makes predictions using the model
\[
y^i = w^Tx^i + w_0.
\]
\blu{Hand in your new function and the updated plot.}
\ans{
Code in \texttt{leastSquares.jl}.
\lstinputlisting{a1/leastSquares.jl}
\centerfig{.5}{yintercept.png}
}
\pagebreak
\item Allowing a non-zero y-intercept improves the prediction substantially, but the model still seems sub-optimal because there are obvious non-linear effects. Write a function called \emph{leastSquaresRBFL2} that implements \emph{least squares using Gaussian radial basis functions (RBFs) and L2-regularization}. \\You should start from the \emph{leastSquares} function and use the same conventions: $n$ refers to the number of training examples, $d$ refers to the number of features, $X$ refers to the data matrix, and $y$ refers to the targets. Use $Z$ as the data matrix after the change of basis, $\sigma$ as the standard deviation in the Gaussian, and $\lambda$ as the regularization parameter. Note that you'll have to add two additional input arguments ($\lambda$ for the regularization parameter and $\sigma$ for the Gaussian RBF variance) compared to the \emph{leastSquares} function. To make your code easier to understand/debug, you may want to define a new function \emph{rbfBasis} which computes the Gaussian RBFs for a given training set, testing set, and $\sigma$ value. \blu{Hand in your function and the plot generated with $\lambda = 1$ and $\sigma = 1$.}\\
Note:  the \emph{distancesSquared} function in \emph{misc.jl} is a vectorized way to quickly compute the squared Euclidean distance between all pairs of rows in two matrices.
\ans{
Code in \texttt{leastSquaresRBFL2.jl}.
\lstinputlisting{a1/leastSquaresRBFL2.jl}
\centerfig{.5}{rbf11.png}
}
\pagebreak
\item Modify the script to split the training data into a ``train'' and ``validation'' set (you can use half the examples for training and half for validation), and use these to select $\lambda$ and $\sigma$. \blu{Hand in your modified script and the plot you obtain with the best values of $\lambda$ and $\sigma$.}
\ans{
Code in \texttt{example\_nonLinear.jl}.
\lstinputlisting[firstline=14,lastline=56]{a1/example_nonLinear.jl}
\centerfig{.5}{bestls.png}
}
\item Consider a model combining the first two parts of this question,
\[
y^i = w^Tx^i + w_0 + v^Tz^i,
\]
where $z^i$ is the Gaussian RBFs and $v$ is a vector of regression weights for those features. Suppose that we first fit $w$ and $w_0$ assuming that $v=0$ (Part 1 of this question), and then we fit $v$ with $w$ and $w_0$ fixed (you could use your code for Part 2 to do this by modifying the $y^i$ values). Why would someone want to do this?
Hint: think about how this model would behave with $x^i = 10$.
\ans{
They might have the prior belief that the trend of the data is going down.
Without the linear regression first, the predicted value of $y_i$ when $x_i=10$ may not trend downward as seen in Q3,
but there may be reason to believe that the trend should continue downward.
}
}
\pagebreak

\subsection{Multi-Class Logistic Regression}


The script \emph{example\_multiClass.jl} loads a multi-class classification dataset and fits a ``one-vs-all'' logistic regression classifier using the \emph{findMin} gradient descent implementation, then reports the validation error and shows a plot of the data/classifier. The performance on the validation set is ok, but could be much better. For example, this classifier never  predicts some of the classes.

Using a one-vs-all classifier hurts performance because the classifiers are fit independently, so there is no attempt to calibrate the columns of the matrix $W$. An alternative to this independent model is to use the softmax probability,
\[
p(y^i | W, x^i) = \frac{\exp(w_{y^i}^\top x^i)}{\sum_{c=1}^k\exp(w_c^\top x^i)}.
\]
Here $c$ is a possible label and $w_{c}$ is column $c$ of $W$. Similarly, $y^i$ is the training label, $w_{y^i}$ is column $y^i$ of $W$. The loss function corresponding to the negative logarithm of the softmax probability is given by
\[
f(W) = \sum_{i=1}^n \left[-w_{y^i}^\top x^i + \log\left(\sum_{c' = 1}^k \exp(w_{c'}^\top x^i)\right)\right].
\]
Make a new function, \emph{softmaxClassifier}, which fits $W$ using the softmax loss from the previous section  instead of fitting $k$ independent classifiers. \blu{Hand in the code and report the validation error}.

Hint: you can use the \emph{derivativeCheck} option when calling \emph{findMin} to help you debug the gradient of the softmax loss. Also, note that the \emph{findMin} function treats the parameters as a vector (you may want to use \emph{reshape} when writing the softmax objective).
\ans{
Code in \texttt{softmaxClassifier.jl}.
\lstinputlisting{a1/softmaxClassifier.jl}
Obtained \texttt{trainError = 0.004, validError = 0.026}.
Referenced \url{https://www.cs.mcgill.ca/~siamak/COMP551/slides/4-logistic-regression_short.pdf}.
\centerfig{.5}{softmax.png}
}
\pagebreak
 
\section{Calculation Questions}

\subsection{Matrix Notation, Quadratics, Convexity, and Gradients}

This question reviews several comptutional tools covered in CPSC 340. You may also find some of the notes on the course webpage useful.

\enum{
\item Consider the function
\[
f(x) = \sum_{i=1}^n\sum_{j=1}^n a_{ij}x_ix_j + \sum_{i=1}^n b_ix_i + c,
\]
where $x$ is a vector of length $n$ with elements $x_i$, $b$ is a vector of length $n$ with elements $b_i$, and $A$ is an $n \times n$ matrix with elements $a_{ij}$ (not necessarilyn symmetric). \blu{Write this function in matrix notation} (so it uses $A$ and $b$, and does not have summations or references to indices $i$.
\\
\[f(\mathbf{x}) = \mathbf{x}^T \mathbf{Ax} + \mathbf{x}^T \mathbf{b} + c\]
\item \blu{Write the gradient of $f$ from the previous question in matrix notation}.
\\
\[\nabla f(\mathbf{x}) = (\mathbf{A+A}^T)\mathbf{x} + \mathbf{b}\]
\item \blu{Give a linear system whose solution gives a minimum value of $f$ in terms of $x$}, in the case where $A$ is symmetric and positive semi-definite (which implies that $f$ is convex).
\\
\[2\mathbf{Ax+b=0}\]
\pagebreak
\item Consider weighted linear regression with an L2-regularizer with regularization strength $1/\sigma^2$,
\[
f(w) = \half \sum_{i=1}^n v^i(y^i - w^Tx^i)^2 + \frac{1}{2\sigma^2}\sum_{j=1}^d w_j^2,
\]
where $\sigma > 0$ and we have a vector $v$ of length $n$ containing the elements $v^i$. The other sybmbols are our standard regression notation. \blu{Write this function in matrix notation}.\\
Note: you and use $X$ as the matrix containing $x^i$ as the rows, $y$ as the vector containing the elements $y^i$, and $V$ as a diagonal matrix containing the vector $v$ along the diagonal.
\\
\[f(\mathbf{w}) = \mathbf{(Xw-y)}^T \mathbf{V(Xw-y)} + \frac{1}{2\sigma^2} \mathbf{w}^T \mathbf{w}\]
\item Assuming that we have $v^i \geq 0$ for all $i$, \blu{show that $f$ from the previous part is convex}.
\[f(\mathbf{w}) = \|\mathbf{V}^{1/2} \mathbf{(Xw-y)}\|^2 + \frac{1}{2\sigma^2} \|\mathbf{w}\|^2\]
Since $f$ is the sum of 2 squared norms, it is convex.
\item Assuming that we have $v^i \geq 0$ for all $i$, \blu{give a linear system whose solution gives a minimum value of $f$ in terms of $w$}.
\\
\[\nabla f(\mathbf{w}) = 2\mathbf{X}^T \mathbf{V(Xw-y)} + \frac{1}{\sigma^2} \mathbf{w}=0\]
\pagebreak
\item If we fit a linear classifer with the exponential loss (used in older boosting algorithms), it would take the form
\[
f(w) = \sum_{i=1}^n \exp(-y^iw^Tx^i).
\]
\blu{Derive the gradient of this loss function.}
\\
\[\nabla f(\mathbf{w}) = -\sum_{i=1}^{n} y^i \mathbf{x}^i \exp{(-y^i \mathbf{w}^T \mathbf{x}^i)}\]
\item The support vector regression objective function is
\[
f(w) = \sum_{i=1}^n \max\{0, |w^Tx^i - y^i|-\epsilon\} + \frac \lambda 2 \norm{w}^2.
\]
where $\epsilon$ is a non-negative hyper-parameter. It is similar to the L1 robust regression loss, but where there is no penalty for being within $\epsilon$ of the prediction (which can reduce overfitting).\blu{ Show that this non-differentiable function is convex}.
\ans{
Norms are convex, so $\norm{w}^2$ is convex. 
$\mathbf{w}^T \mathbf{x}^i$ is convex so summing with a constant, taking the L1-norm and applying the max function with another constant will be convex as well.
Summing the above convex functions results in $f$ being convex as well.
}
}
\pagebreak



\subsection{MAP Estimation}


In 340, we showed that under the assumptions of a Gaussian likelihood and Gaussian prior,
\[
y^i \sim \mathcal{N}(w^\top x^i,1), \quad w_j \sim \mathcal{N}\left(0,\frac{1}{\lambda}\right),
\]
that the MAP estimate is equivalent to solving the L2-regularized least squares problem
\[
f(w) = \frac{1}{2}\sum_{i=1}^n (w^\top x^i - y^i)^2 + \frac \lambda 2 \sum_{j=1}^d w_j^2,
\]
in the ``loss plus regularizer'' framework.
For each of the alternate assumptions below, write it in the ``loss plus regularizer'' framework (simplifying as much as possible):

\blu{\enum{
\item Gaussian likelihood with a separate variance $\sigma_i^2$ for each training example, and Laplace prior with a separate variance $1/\lambda_j$ for each variable,
\[
y^i \sim \mathcal{N}(w^Tx^i,\sigma_i^2), \quad w_j \sim \mathcal{L}\left(0,\frac{1}{\lambda_j}\right).
\]
\[p(y^i|x^i,w,\sigma_i) \propto \frac{1}{\sigma_i} \exp{-\frac{(w^T x^i - y^i)^2}{2\sigma_i^2}}\]
\[p(w_j|\lambda_j) \propto \lambda_j \exp{-\lambda_j |w_j|}\]
Equivalent to solving
\[f(w)=\sum_{i=1}^n \left(\frac{(w^T x^i - y^i)^2}{2\sigma_i^2} + \log{\sigma_i}\right)
+ \sum_{j=1}^d \left(\lambda_j |w_j| - \log{\lambda_j}\right)\]
\item Robust student-$t$ likelihood and Gaussian prior centered at $u$.
\[
p(y^i | x^i, w) = \frac{1}{\sqrt{\nu}B\left(\frac 1 2,\frac \nu 2\right)}\left(1 + \frac{(w^Tx^i - y^i)^2}{\nu}\right)^{-\frac{\nu + 1}{2}}, \quad w_j \sim \mathcal{N}\left(u_j,\frac{1}{\lambda}\right),
\]
where $u$ is $d \times 1$, $B$ is the ``Beta" function, and the parameter $\nu$ is called the ``degrees of freedom".\footnote{This likelihood is more robust than the Laplace likelihood, but leads to a non-convex objective.}
\[f(w)=\frac{\nu+1}{2}\sum_{i=1}^n\log{\left(1+\frac{(w^T x^i - y^i)^2}{\nu}\right)} 
+ \frac{\lambda}{2}\sum_{j=1}^d (w_j-u_j)^2\]
\item We use a Poisson-distributed likelihood (for the case where $y_i$ represents counts), and we use a uniform prior for some constant $\kappa$,
\[
p(y^i | x^i, w) = \frac{\exp(y^iw^\top x^i)\exp(-\exp(w^\top x^i))}{y^i!}, \quad p(w_j) \propto \kappa.
\]
(This prior is 	``improper'' since $w\in\R^d$ but it doesn't integrate to 1 over this domain.)
\[f(w)=\sum_{i=1}^n \left(\exp{(w^T x^i)} - y^i w^T x^i + \sum_{k=1}^{y^i} \log k \right)\]
}
}


\pagebreak

\subsection{Machine Learning Model Memory and Time Complexities}


Answer the following questions using big-O notation, and a brief explanation.
Your answers may involve $n$, $d$, and perhaps additional quantities defined in the question. 
As an example, the (linear) least squares model has $O(d)$ parameters, requires $O(d)$ time for prediction, and requires $O(nd^2 + d^3)$ time to train.\footnote{In this course, we assume matrix operations have the ``textbook'' cost where the operations are implemented in a straightforward way with ``for'' loops. For example, we'll assume that multiplying two $n \times n$ matrices or computing a matrix inverse simply costs $O(n^3)$, rather than the $O(n^\omega)$ where $\omega$ is closer to $2$ as discussed in CS algorithm courses.}

\blu{
\enum{
\item What is the storage space required for the $k$-means clustering algorithm?
\ans{
\[O((n+k)d)\]
$O(nd)$ for the data, and $O(kd)$ for the cluster locations. Referenced StackOverflow: \url{https://stackoverflow.com/a/25362871}
}
\item What is the cost of clustering $t$ examples using an already-trained k-means model?
\ans{
\[O(tdk)\]
$O(dk)$ to compute distance to $k$ clusters for each example, $t$ times.
}
\item What is the training time for linear regression with L2-regularization?
\ans{
\[O(nd^2 + d^3)\]
Same as the least squares.
}
\item What is the prediction cost for linear regression with L2-regularization on $t$ test examples?
\ans{
\[O(td)\]
$O(d)$ for one example, $t$ times.
}
\pagebreak
\item What is the storage space required for linear regression with Gaussian RBFs as features?
\ans{
\[O(n^2 + nd)\]
$O(nd)$ for the training examples, $O(n^2)$ for the new basis.
}
\item What is the prediction time for linear regression with Gaussian RBFs as features on $t$ test examples? You can use $\sigma^2$ as the variance of the Gaussian RBFs.
\ans{
\[O(nt)\]
$O(n)$ for each example, $t$ times.
}
\item What is the cost of evaluating the support vector regression objective function?
\ans{
\[O(nd)\]
Multiplying $O(d)$ numbers together for $w^T x_i, n$ times.
}
\item What is the cost of trying to minimize the exponential loss regression loss by running $t$ iterations of  gradient descent?
\ans{
\[O(ndt)\]
$O(nd)$ to calculate the loss for a single iteration, $t$ times.
}
\item What is the cost of trying to minimize the exponential loss  by running $t$ iterations of stochastic gradient descent?
\ans{
\[O(dt)\]
$O(d)$ to calculate loss for a single iteration with 1 training example, $t$ times.
}
}
}
Hint: many people got very-low grades on this question last year. If you are not sure how to answer questions like this, get help!


 
\end{document}